\documentclass[11pt]{article}
\usepackage[T1]{fontenc}
\usepackage{lmodern,amsmath,amssymb,amsthm,mathtools,microtype}
\usepackage[margin=1in]{geometry}
\usepackage[colorlinks=true,linkcolor=blue,citecolor=blue,urlcolor=blue]{hyperref}
\newtheorem{theorem}{Theorem}[section]
\newtheorem{proposition}[theorem]{Proposition}
\newtheorem{lemma}[theorem]{Lemma}
\newtheorem{corollary}[theorem]{Corollary}
\theoremstyle{remark}
\newtheorem{remark}[theorem]{Remark}
\newcommand{\D}{\mathcal D}
\newcommand{\ddc}{dd^c}
\newcommand{\supp}{\operatorname{supp}}
\newcommand{\osc}{\operatorname{osc}}
\title{Intrinsic current methods for the K\"ahler surface regularization problem}
\author{Research note}
\date{October 29, 2026}
\begin{document}
\maketitle
\begin{abstract}
We investigate whether a continuous plurisubharmonic exhaustion on an arbitrary noncompact K\"ahler surface can be replaced by a smooth one. The general implication is not proved here, and no counterexample is constructed. We establish intrinsic averaging and $L^2$ closedness statements, give a complete infinite-volume metric construction, and prove a singular-current closedness result under bounded curvature. A separate two-current argument treats continuously exhausted domains in compact K\"ahler surfaces without any divisorial-boundary hypothesis. We formulate the remaining separation problem by smooth psh hulls and prove a bounded-mass duality criterion that exposes the missing closed-range step. We also prove a conditional regularization theorem, without compactification, when a complete bounded-curvature K\"ahler metric has a positive Green kernel tending to zero at infinity.
\end{abstract}

\section{The question and reductions}
Throughout, $X$ is a connected noncompact complex surface and an exhaustion means a real-valued function with compact sublevel sets. The target is
\[
 X\text{ K\"ahler},\quad \rho\in C^0(X)\cap\operatorname{PSH}(X)
 \text{ exhaustive}
 \quad\Longrightarrow\quad
 \operatorname{PSH}(X)\cap C^\infty(X)
 \text{ contains an exhaustion}.
\]
We set $\D=i\partial\bar\partial$; positive normalization factors in $dd^c$ do not affect the assertions. The existence of some smooth exhaustion is requested, not smooth approximation of every given psh function. The compact-divisor-complement proof does not imply this assertion.

Two established reductions do not require a compactification. If $\mathcal O(X)\ne\mathbb C$, V\^aj\^aitu's holomorphic-convexity theorem \cite[Theorem~1.1]{Vajaitu} applies, and a smooth strictly psh exhaustion of the Stein Remmert reduction pulls back to a smooth psh exhaustion of $X$. This is the safe citation for the continuous-exhaustion version: Ohsawa's original standing definition of weak completeness used a smooth exhaustion.

If $X$ has exactly one end, the Bochner--Hartogs dichotomy \cite[Theorem~2.5, case~(i)]{NR} says that either $H_c^1(X,\mathcal O)=0$, or there is a proper holomorphic map from $X$ onto a Riemann surface. The latter surface is noncompact, because the map is proper and $X$ is noncompact; pulling back its smooth strictly subharmonic exhaustion solves our problem. Thus a one-ended counterexample must have
\[
 \mathcal O(X)=\mathbb C,\qquad H_c^1(X,\mathcal O)=0.
\]
Neither condition is itself an obstruction to smooth exhaustion.

\section{Intrinsic averaging}
\begin{proposition}\label{prop:averaging}
Let $X$ be a complex surface with a continuous psh exhaustion $\rho$ and let $\omega$ be any smooth Hermitian form. If $X$ is not a modification of a Stein space, there is $t_0$ such that for every $t_0<a<b$ there is a positive $(1,1)$-current $T$ satisfying
\[
 \D T=0,\quad \supp T\subset\rho^{-1}([a,b]),\quad
 \int_XT\wedge\omega=1,\quad \nu(T,x)=0\quad(x\in X).
\]
In particular, compactification and K\"ahlerness are unnecessary for this step.
\end{proposition}
\begin{proof}
Let $\Sigma_X^0$ be the continuous minimal kernel. By \cite[Lemma~2.1]{MS}, failure to be a Stein modification implies
$(t_0,\infty)\subset\rho(\Sigma_X^0)$ for some $t_0$.
For $t>t_0$, the compact level $K_t=\{\rho=t\}$ cannot have a smooth strictly psh neighborhood function $v$. Indeed, normalize $0<v<1$ on a narrow compact band $|\rho-t|\le\delta$, choose $0<\eta<1$ and $0<\varepsilon<\eta\delta/2$, and replace
$g(\rho)$, where $g(s)=s+\eta|s-t|$, on that band by
\[
 \max\{g(\rho),\,\rho+\varepsilon v\}.
\]
The first branch dominates near the edges and the second near $K_t$.
This produces a continuous psh exhaustion strongly psh near the entire level, contradicting $K_t\cap\Sigma_X^0\ne\varnothing$.

Separate, in the Banach space of continuous real $(1,1)$-forms on $K_t$, the open positive-definite cone from
$\{\D f|_{K_t}:f\in C_c^\infty(X,\mathbb R)\}$.
Hahn--Banach gives a positive order-zero current $T_t$ supported on $K_t$, annihilating all $\D f$. Its positive Hermitian mass may be normalized to one.

For each $n$, choose the $n$ equally spaced midpoint levels in $[a,b]$ and average their currents:
\[
 t_{n,j}=a+(j-\tfrac12)(b-a)/n,\qquad
 T_n=\frac1n\sum_{j=1}^nT_{t_{n,j}}.
\]
Their supports lie in one compact set, so a subsequence converges weakly to a positive $\D$-closed current of mass one. No measurable selection is needed.

In a coordinate ball, let $\beta=\D|z|^2$ and fix $R$ with the closed $2R$ ball in the chart. The local Lelong--Jensen formula for positive $\D$-closed currents \cite[Appendix~B]{DNS} gives
\[
 r^{-2}\int_{B(x,r)}T_t\wedge\beta
 \le R^{-2}\int_{B(x,R)}T_t\wedge\beta\le C_R.
\]
Only grid levels in
$[\inf_{B(x,r)}\rho,\sup_{B(x,r)}\rho]$ contribute. Thus
\[
 r^{-2}\int_{B(x,r)}T_n\wedge\beta
 \le C_R\left(\frac{\osc_{B(x,r)}\rho}{b-a}+\frac2n\right).
\]
Use a nonnegative compactly supported cutoff times $\beta$ to regard current convergence as weak convergence of positive measures. The open-set Portmanteau inequality gives
\[
 r^{-2}\int_{B(x,r)}T\wedge\beta
 \le\frac{C_R}{b-a}\osc_{B(x,r)}\rho\longrightarrow0.
\]
The same convergent subsequence works for every fixed ball. This proves zero Lelong numbers everywhere.
\end{proof}

\section{Complete metrics and an intrinsic intersection inequality}

% Begin complete-infinite-volume
\begin{proposition}\label{prop:complete-infinite-volume}
Let $X$ be a connected noncompact K\"ahler manifold of complex
dimension $N\geq1$. If $X$ admits a continuous plurisubharmonic
exhaustion, then $X$ admits a smooth complete K\"ahler metric of
infinite volume. More precisely, for every initial K\"ahler form
$\omega$, such a form can be chosen as
\[
 \widehat\omega=3\omega+i\partial\bar\partial(\zeta^2),
\]
where $\zeta$ is a positive smooth exhaustion. No smooth
plurisubharmonic exhaustion is assumed or produced by this argument.
\end{proposition}

\begin{proof}
Write $\D=i\partial\bar\partial$. We use the following precise
form of Demailly's approximation theorem
\cite[Theorem~9.4, p.~506]{Demailly82}. There are positive continuous
exhaustions $m,M$ on $X$, with $m<M$, such that for every positive
continuous function $\varepsilon$ on $X$ there exists a smooth
real function $\psi$ satisfying
\begin{equation}\label{eq:arbitrary-error-approximation}
 m\leq\psi\leq M,
 \qquad \D\psi\geq-\varepsilon\omega.
\end{equation}
The functions $m$ and $M$ are fixed before $\varepsilon$ is chosen.
In that reference weak pseudoconvexity is defined by the existence
of a plurisubharmonic exhaustion, without a smoothness requirement
(Definition~1.1, p.~461). Thus the theorem applies to the hypothesis
of the proposition.

Choose a smooth relatively compact connected domain $\Omega_1$
with nonempty interior. Inductively construct smooth relatively
compact connected domains $\Omega_j$, with $K_j=\overline{\Omega_j}$,
and positive numbers $a_j$ so that
\begin{equation}\label{eq:condenser-boundaries}
 K_j\Subset\Omega_{j+1},\qquad
 \bigcup_j\Omega_j=X,\qquad
 M<a_j\text{ on }K_j,\qquad
 m>a_j+1\text{ on }\partial\Omega_{j+1}.
\end{equation}
We may also require $a_{j+1}>a_j+2$. To do this, after choosing
$K_j$, take $a_j>\sup_{K_j}M$, increasing it to meet the previous
numerical requirement. Choose $\Omega_{j+1}$ to contain $K_j$,
the compact set $\{m\leq a_j+1\}$, and the next member of any
fixed compact exhaustion of $X$. These compact sets can be
contained in a smooth relatively compact connected domain.

Put $A_j=\Omega_{j+1}\setminus K_j$. For a real Sobolev function
$v$ on a relatively compact set $A$, use the Dirichlet energy
\[
 \mathcal E_A(v)
 =\int_A i\partial v\wedge\bar\partial v\wedge\omega^{N-1}.
\]
Define the condenser capacity
\[
 c_j=\inf\left\{\mathcal E_{A_j}(v):
 v\in W^{1,2}(A_j),\quad
 \operatorname{Tr}_{\partial\Omega_j}v=0,\quad
 \operatorname{Tr}_{\partial\Omega_{j+1}}v=1\right\}.
\]
Here all traces are taken from $A_j$. We have $0<c_j<\infty$.
Finiteness follows by choosing a smooth function equal to zero
near $K_j$ and one near $\partial\Omega_{j+1}$. For positivity,
extend an admissible $v$ by zero across $\Omega_j$. Then $1-v$
belongs to $W^{1,2}_0(\Omega_{j+1})$ and equals one on $\Omega_j$.
The Poincar\'e inequality on the bounded smooth connected domain
$\Omega_{j+1}$ gives
\[
 0<\operatorname{Vol}_{\omega}(\Omega_j)
 \leq\int_{\Omega_{j+1}}|1-v|^2\,dV_\omega
 \leq C_j\int_{A_j}|dv|_\omega^2\,dV_\omega
 =C'_j\mathcal E_{A_j}(v).
\]
The last equality is the usual fixed positive dimensional
normalization between these two Dirichlet energies. Taking the
infimum proves strict positivity, even if $A_j$ is disconnected.

Choose a smooth increasing convex function
$\chi:[0,\infty)\to[1,\infty)$, with $\chi'\geq1$, such that
\begin{equation}\label{eq:large-increments}
 \delta_j:=\chi(a_j+1)-\chi(a_j)\geq c_j^{-1/2}
 \quad\text{for every }j.
\end{equation}
Such a function exists: choose a smooth nondecreasing positive
derivative, at least one everywhere, whose value on each
$[a_j,a_j+1]$ is at least $c_j^{-1/2}$, and integrate it.

Only now choose the error function
\[
 \varepsilon(x)=
 \frac{1}{\chi(M(x))\chi'(M(x))}>0
\]
and obtain $\psi$ from \eqref{eq:arbitrary-error-approximation}.
Set $\zeta=\chi(\psi)$. Since $\chi$ is increasing and convex,
\begin{align*}
 \D\zeta
 &=\chi'(\psi)\D\psi
   +\chi''(\psi)i\partial\psi\wedge\bar\partial\psi\\
 &\geq-\chi'(\psi)\varepsilon\omega
 \geq-\frac{\omega}{\chi(M)}.
\end{align*}
We also have $0<\zeta\leq\chi(M)$. Therefore
\begin{equation}\label{eq:complete-metric-bound}
 \widehat\omega
 :=3\omega+\D(\zeta^2)
 =3\omega+2\zeta\D\zeta
       +2i\partial\zeta\wedge\bar\partial\zeta
 \geq\omega+2i\partial\zeta\wedge\bar\partial\zeta.
\end{equation}
In particular, $\widehat\omega$ is a smooth positive closed
$(1,1)$-form. Notice that the bound $\psi\leq M$ is essential
in both estimates; simply composing an arbitrary almost
plurisubharmonic function with $\chi$ would not suffice.

Because $\zeta\geq\chi(m)$, the function $\zeta$ is proper.
Inequality~\eqref{eq:complete-metric-bound} bounds $|d\zeta|$
with respect to the Riemannian metric of $\widehat\omega$ by a
fixed constant. Consequently every curve escaping all compact
subsets has infinite length. Equivalently, any Cauchy sequence for
that metric has bounded $\zeta$ and hence lies in a compact
subset, where it converges. Thus $\widehat\omega$ is complete.

It remains to prove that its volume is infinite. From
\eqref{eq:condenser-boundaries} and $m\leq\psi\leq M$ we have
\[
 \zeta\leq\chi(a_j)\text{ on }\partial\Omega_j,
 \qquad
 \zeta\geq\chi(a_j+1)\text{ on }\partial\Omega_{j+1}.
\]
The Lipschitz truncation
\[
 v_j=\min\left\{1,\max\left\{0,
       \frac{\zeta-\chi(a_j)}{\delta_j}\right\}\right\}
\]
is therefore admissible for $c_j$. The Sobolev chain rule gives
\[
 c_j\leq\mathcal E_{A_j}(v_j)
 \leq\delta_j^{-2}\mathcal E_{A_j}(\zeta),
 \qquad
 \mathcal E_{A_j}(\zeta)\geq c_j\delta_j^2\geq1.
\]
The annuli $A_j$ are pairwise disjoint. Hence
\[
 \int_X i\partial\zeta\wedge\bar\partial\zeta
                      \wedge\omega^{N-1}=\infty.
\]
Finally, monotonicity of the determinant for positive Hermitian
forms and the rank-one identity
$(i\partial\zeta\wedge\bar\partial\zeta)^2=0$ imply
\[
 \widehat\omega^N
 \geq\bigl(\omega+2i\partial\zeta
                         \wedge\bar\partial\zeta\bigr)^N
 =\omega^N+2N i\partial\zeta\wedge\bar\partial\zeta
                                      \wedge\omega^{N-1}.
\]
Its integral, and therefore the volume of $\widehat\omega$,
is infinite.
\end{proof}

\begin{remark}
Demailly's Theorem~1.3 already gives completeness, using the
same square-potential construction with a suitably small error.
The condenser argument above is an additional deduction ensuring
infinite volume. It does not supply curvature bounds, bounded
geometry, a Green function, or estimates for regularizing
singular currents. In particular it does not resolve the gap
between measure-valued positive $\partial\bar\partial$-closed
currents and an $L^2$ closedness theorem.
\end{remark}


% End complete-infinite-volume


% Begin intrinsic-L2
\begin{theorem}\label{thm:L2-closedness}
Let $(X,\omega)$ be a connected complete K\"ahler surface of infinite
volume. Let $T$ be a real $(1,1)$-current whose coefficients are globally
in $L^2(X,\omega)$. Suppose $T\geq0$ and $i\partial\bar\partial T=0$
in the sense of distributions. Then $dT=0$. Moreover, $T\wedge T=0$
almost everywhere, where this wedge product is the $L^1$ form defined
by the $L^2$ coefficients.
\end{theorem}

\begin{proof}
Set $\D=i\partial\bar\partial$, let $\Lambda$ be contraction with
$\omega$, and put $dV=\omega^2/2$. These choices only fix constants;
the conclusion is unchanged if $dd^c$ is any positive multiple of
$\D$. For every real $(1,1)$-form $\gamma$ on a K\"ahler surface,
the pointwise linear-algebra identity is
\begin{equation}\label{eq:linear-algebra}
 \gamma\wedge\gamma
   =\bigl((\Lambda\gamma)^2-|\gamma|^2\bigr)dV.
\end{equation}
For example, in a unitary frame this is the identity
$2\lambda_1\lambda_2=(\lambda_1+\lambda_2)^2-
(\lambda_1^2+\lambda_2^2)$.
In particular, a real primitive $(1,1)$-form satisfies
$\gamma\wedge\gamma=-|\gamma|^2dV$.

We first justify a density statement for the scalar operator
$L=\Lambda\D$ with domain $C_c^\infty(X,\mathbb R)$.
Up to a fixed nonzero real factor and sign, this is the scalar
Laplace--Beltrami operator. If an $L^2$ function $h$ is orthogonal to the range of $L$, it is distributionally harmonic and hence smooth by elliptic regularity. Complete-metric cutoffs $\chi_R$ with $|d\chi_R|\leq C/R$ and the Caccioppoli inequality give
\[
 \int_X\chi_R^2|dh|^2\,dV
 \leq4\int_X|h|^2|d\chi_R|^2\,dV\longrightarrow0.
\]
Thus $h$ is constant; connectedness and infinite volume force that constant to vanish.
Consequently
\[
 \overline{L(C_c^\infty(X,\mathbb R))}^{\,L^2}
  = (\ker L^*)^\perp=L^2(X,\mathbb R).
\]
Only density of the range is used; no closed-range assertion and no
bounded Green operator are needed. The same vanishing of the scalar
$L^2$ kernel is used in \cite[Lemma~3.1.1]{Dingoyan}.

Put $f=\Lambda T\in L^2(X,\mathbb R)$. Choose real
$u_n\in C_c^\infty(X)$ such that $Lu_n\longrightarrow f$ in $L^2$.
For every real $v\in C_c^\infty(X)$, Stokes' theorem and
\eqref{eq:linear-algebra} give
\[
  \int_X(\D v)^2=0,
  \qquad
  \|\D v\|_{L^2}^2=\|\Lambda\D v\|_{L^2}^2.
\]
The first equality follows because $\D v$ is an exact closed form
with compact support. Applying the second equality to $u_n-u_m$
shows that $\eta_n=\D u_n$ is Cauchy in $L^2$. Let its limit be
$\eta$. Then $\eta$ is a real closed $(1,1)$-current,
$\Lambda\eta=f$, and
\[
 \int_X\eta^2=\lim_n\int_X\eta_n^2=0.
\]
Here and below convergence of the integrals follows from the
continuity of the wedge pairing on $L^2\times L^2$.

The distributional equation $\D T=0$, tested against $u_n$, says
$\int_XT\wedge\eta_n=0$. Hence $\int_XT\wedge\eta=0$.
The form $\beta=T-\eta$ is primitive, so
\[
 \int_X T^2
  =\int_X(T-\eta)^2
  =-\|\beta\|_{L^2}^2\leq0.
\]
On the other hand, positivity of $T$ implies $T^2\geq0$
almost everywhere. Therefore $\beta=0$ and $T=\eta$ is closed.
The same identity gives $\int_XT^2=0$, so $T^2=0$ almost everywhere.
\end{proof}

\begin{remark}
Compact support is not needed in Theorem~\ref{thm:L2-closedness}. Proposition~\ref{prop:complete-infinite-volume} supplies a metric with its completeness and volume hypotheses on the surfaces under investigation. The remaining issue is that the currents from averaging need not have $L^2$ coefficients.
\end{remark}

\begin{remark}[Averaging can remain singular]
Let $E$ be a compact elliptic curve, let $X=E\times\mathbb C$, and
give $X$ a product complete K\"ahler metric with
$\int_E\omega_E=1$. The metric has infinite volume. For
$0<a<b$, take the smooth psh exhaustion $\rho(e,z)=|z|^2$ and set
\[
 T=\frac1{b-a}\int_a^b [E\times\{\sqrt t\}]\,dt.
\]
This is precisely an average of mass-one positive closed currents
supported on the levels $\rho=t$. It is positive and closed, has
compact support, and has zero Lelong number at every point:
its local trace mass is $O(r^3)$, including at the two endpoints
of the interval. But its trace measure is a nonzero measure supported
on the real three-dimensional subset
$E\times[\sqrt a,\sqrt b]$, which has zero four-dimensional volume.
Thus $T$ has no $L^2$ coefficients. Even a smooth exhaustion and the
vanishing of every Lelong number do not bridge this gap.
\end{remark}




% End intrinsic-L2


\section{Singular currents under bounded curvature}

% Begin bounded-geometry
\begin{proposition}\label{prop:bounded-geometry-closedness}
Let $(X,\omega)$ be a connected complete K\"ahler surface of infinite
volume whose Riemannian curvature tensor is uniformly bounded.
Let $T$ be a positive $(1,1)$-current with compact support
such that
\[
 i\partial\bar\partial T=0,
 \qquad \nu(T,x)=0\quad\hbox{for every }x\in X.
\]
Then $dT=0$.
\end{proposition}

We use $\D=i\partial\bar\partial$, $dV=\omega^2/2$, and contraction
$\Lambda$ with $\omega$. Any positive normalization of $dd^c$ gives
the same assertions. The scalar nonnegative Laplacian is denoted by
$\Delta_0$, its heat semigroup by $P_t=e^{-t\Delta_0}$, and the Hodge
heat semigroup on real $(1,1)$-forms by $H_t=e^{-t\Delta_2}$.

\begin{lemma}[$L^2$ primitive-defect identity]\label{lem:L2-defect}
On a connected complete infinite-volume K\"ahler surface, let $A$ be
a real $L^2$ $(1,1)$-form with $\D A=0$ distributionally. There is a
real closed $L^2$ $(1,1)$-form $\eta$ such that
\[
 \Lambda\eta=\Lambda A,\qquad
 \int_X\eta^2=\int_X A\wedge\eta=0,
 \qquad
 \int_X A^2=-\|A-\eta\|_{L^2}^2.
\]
No positivity assumption on $A$ is required.
\end{lemma}

\begin{proof}
The construction of $\eta$ in the proof of Theorem~\ref{thm:L2-closedness} used only the real $L^2$ and $\D$-closed hypotheses. Positivity was used only after obtaining the displayed identity. Repeating that construction with $A$ gives all the assertions.
\end{proof}

\begin{proof}[Proof of Proposition~\ref{prop:bounded-geometry-closedness}]
Let $\mu=T\wedge\omega$ be the trace measure. It is finite and
compactly supported. Positivity implies that the tensor variation
of $T$ is bounded by a fixed multiple of $\mu$.

\medskip\noindent
\emph{Uniform local mass decay.}
We first show
\begin{equation}\label{eq:uniform-decay}
 \varepsilon(r):=\sup_{x\in X}\frac{\mu(B_g(x,r))}{r^2}
       \longrightarrow0\quad(r\downarrow0).
\end{equation}
In a holomorphic coordinate chart put $\beta=\D|z|^2$. On a smaller
compact set of centers, for every sufficiently small fixed $r>0$,
the function
\[
 x\longmapsto r^{-2}
       \int_{\overline{B_{\rm eucl}(x,r)}}T\wedge\beta
\]
is upper semicontinuous. Indeed, the closed balls vary continuously
and their measures obey the closed-set upper-semicontinuity
inequality. The Lelong--Jensen formula for positive $dd^c$-closed
currents makes these normalized masses nondecreasing in $r$, with
limit zero as $r\downarrow0$; see
\cite[Appendix~B, (B.1)--(B.2)]{DNS}.
Given $\epsilon>0$, the open sets on which a fixed-radius normalized
mass is less than $\epsilon$ cover the compact set of centers as the
radius varies toward zero. A finite subcover and monotonicity give
uniform convergence to zero. This is the upper-semicontinuous
version of the elementary Dini argument; closed balls, rather than
open balls, are used at this step.

A finite chart cover of a compact neighborhood of $\supp T$, and
uniform comparisons of $g,\omega$ with their coordinate Euclidean
counterparts there, give \eqref{eq:uniform-decay}. For small $r$,
balls contributing positive mass have centers in this fixed compact
neighborhood. The same notation $\varepsilon(r)$ may be used for
all $r>0$; it is bounded, because of \eqref{eq:uniform-decay}, local
mass bounds, and $\mu(X)<\infty$.

\medskip\noindent
\emph{Scalar heat bounds.}
Set $q_t=P_t\mu$ and $K=\supp\mu$. For all sufficiently small $t>0$,
all $x\in X$, and all $y\in K$, the scalar heat kernel satisfies
\begin{equation}\label{eq:Gaussian}
 0\leq p_t(x,y)\leq C t^{-2}
       \exp\bigl(-d_g(x,y)^2/(Ct)\bigr).
\end{equation}
Indeed, the scalar Gaussian bound recorded in
\cite[Introduction, before Proposition~1.1]{BDG}, applied using
symmetry, has the factor $\operatorname{Vol}(B_g(y,\sqrt t))^{-1}$.
Bounded curvature gives the required Ricci lower bound. Since the
source point $y$ lies in the fixed compact set $K$, smoothness of
the metric gives $\operatorname{Vol}(B_g(y,\sqrt t))\ge c_Kt^2$
uniformly there. No global injectivity-radius bound is required.
Dyadic annuli of
radii $2^k\sqrt t$, together with \eqref{eq:uniform-decay}, give
\[
 t\|q_t\|_\infty
 \leq C\sum_{k\geq0}4^k e^{-c4^k}
                  \varepsilon(2^{k+1}\sqrt t)=o(1).
\]
Adjusting the first finitely many constants absorbs the innermost
ball. Dominated convergence applies because $\varepsilon$ is bounded
and tends to zero at zero. Sub-Markovianity gives
$\|q_t\|_1\leq\mu(X)$, and therefore
\begin{equation}\label{eq:heat-L2}
 t\|q_t\|_2^2
 \leq t\|q_t\|_\infty\|q_t\|_1=o(1).
\end{equation}

\medskip\noindent
\emph{Heat regularization and its positivity defect.}
The Weitzenb\"ock formula on real $(1,1)$-forms is
$\Delta_2=\nabla^*\nabla+\mathcal R$, with
$\|\mathcal R\|_\infty\leq C$ under the curvature bound.
Let $Q_t=e^{-t\nabla^*\nabla}$. The metric connection preserves
the real $(1,1)$ bundle and its positive Hermitian cone. Its heat
semigroup is an average of parallel transports, so $Q_t$ preserves
this cone and satisfies $|Q_t B|\leq P_t|B|$.
The Ricci lower bound ensures stochastic completeness, as required
for the global probabilistic representation.
The parallel-transport representation and covariant Schr\"odinger
semigroup domination are given in \cite[Lemmas~3.6--3.7]{BDG}
(the connection-heat case has zero potential); in particular,
\[
 |H_t B|\leq e^{Ct}P_t|B|.
\]
These assertions extend to finite compactly supported
measure-valued initial forms, by kernel domination or positive
smooth approximation in local unitary frames.

Define $A_t=H_tT$. It is smooth, real, and of type $(1,1)$, converges
to $T$ as a current as $t\downarrow0$, and satisfies
\[
 |A_t|\leq Ce^{Ct}q_t.
\]
Duhamel's formula and the scalar semigroup identity yield
\begin{align*}
 |A_t-Q_tT|
 &\leq \int_0^t P_{t-s}(|\mathcal R A_s|)\,ds\\
 &\leq C\int_0^t e^{Cs}P_{t-s}(P_s\mu)\,ds
 \leq Cte^{Ct}q_t.
\end{align*}
Consequently, for $0<t\leq1$,
\begin{equation}\label{eq:positivity-defect}
 |A_t|\leq Cq_t,
 \qquad A_t\geq-Ctq_t\,\omega.
\end{equation}
Local parabolic regularity supplies smoothness without uniform
bounds on derivatives of curvature. The Gaussian estimates and
$q_t\in L^1\cap L^\infty$ imply
$A_t\in L^2$.

The complete K\"ahler identities give
$\Delta_d=2\Delta_\partial=2\Delta_{\bar\partial}$ and the
commutation of Hodge heat with $\partial$, $\bar\partial$, and
$\D$; see \cite[\S12]{DemaillyHodge}.
These are identities for the complete self-adjoint realizations,
initially on smooth compactly supported test forms and then by
closure. Hodge heat also commutes with the Hodge star, so its
self-adjointness gives symmetry for the wedge pairing.
The identities apply to compactly supported currents by transposition.
More explicitly, for $\phi\in C_c^\infty(X)$,
\[
 \langle A_t,\D\phi\rangle
 =\langle T,H_t\D\phi\rangle
 =\langle T,\D P_t\phi\rangle=0.
\]
In the last equality, $P_t\phi$ may be multiplied by a cutoff equal
to one on a neighborhood of $\supp T$, so the distributional
equation $\D T=0$ applies. Thus $\D A_t=0$.

\medskip\noindent
\emph{Vanishing of the primitive defect.}
At each point, let $\lambda_1,\lambda_2$ be the Hermitian eigenvalues
of $A_t$ relative to $\omega$. By \eqref{eq:positivity-defect},
$|\lambda_j|\leq Cq_t$ and $\lambda_j\geq-Ctq_t$. Hence
\[
 A_t^2=2\lambda_1\lambda_2\,dV
             \geq-Ctq_t^2\,dV.
\]
Apply Lemma~\ref{lem:L2-defect} to $A_t$, writing
$A_t=\eta_t+\beta_t$ with $d\eta_t=0$ and
\[
 -\|\beta_t\|_2^2=\int_X A_t^2
          \geq-Ct\|q_t\|_2^2=o(1).
\]
Therefore $\|\beta_t\|_2\to0$. Since $dA_t=d\beta_t$, testing
against a compactly supported smooth form shows $dA_t\to0$ as a
distribution. But $A_t\to T$ as currents, and consequently $dT=0$.
\end{proof}

\begin{remark}[Scope of this deduction]
The argument establishes closedness under the displayed metric
hypotheses. It neither constructs such a bounded-curvature complete
K\"ahler metric from a continuous psh exhaustion nor constructs a
smooth psh exhaustion from the resulting closed current. A further
argument that produces functions separating whole exhaustion levels
would still be required. The proof does not assume a noncompact
$dd^c$ lemma, a bounded Green operator, closed range of the scalar
Laplacian, or an extension of the compact DNS cohomology decomposition.
\end{remark}


% End bounded-geometry


\section{A further conclusion with a compact ambient surface}
The following deduction uses the full DNS mixed density formula, not a claim that $\ddc$-closed currents are automatically closed.

\begin{proposition}\label{prop:disjoint}
Let $(S,\omega)$ be a compact K\"ahler surface. Suppose $T_1,T_2$ are positive $\ddc$-closed $(1,1)$-currents of $\omega$-mass one, with disjoint supports and zero Lelong numbers everywhere. Then
\[
 dT_1=dT_2=0,\qquad [T_1]=[T_2],\qquad [T_1]^2=0.
\]
Their common class is nef.
\end{proposition}
\begin{proof}
Use the decomposition of \cite[\S2.1 and Lemma~B.2]{DNS}
\[
 T_j=\Omega_j+\partial R_j+\bar\partial\overline{R_j}
       +i\partial\bar\partial u_j,
 \qquad \alpha_j=[\Omega_j].
\]
Write $U_j=\bar\partial R_j$, an $L^2$ $(0,2)$-form, and use the positive Hermitian pairing
$\langle U,V\rangle=\int_SU\wedge\overline V$.
Theorem~2.2 of \cite{DNS} associates nonnegative self-density masses $m_{jj}$ and gives
\[
 \alpha_j^2=m_{jj}+2\|U_j\|^2.
\]
The mixed tangent current is zero because the supports are disjoint. Its mass formula is therefore
\[
 \alpha_1\cdot\alpha_2=2\operatorname{Re}\langle U_1,U_2\rangle.
\]
All mass-formula hypotheses hold, since both positive-Lelong-number loci are empty. Subtraction yields
\[
 (\alpha_1-\alpha_2)^2
   =m_{11}+m_{22}+2\|U_1-U_2\|^2\ge0.
\]
Equal masses make $\alpha_1-\alpha_2$ primitive. Hodge index makes its square nonpositive, with equality only for the zero class. Hence
$\alpha_1=\alpha_2$, $m_{11}=m_{22}=0$, and $U_1=U_2$.
Differentiating the decompositions now gives $dT_1=dT_2$.
These two distributions have disjoint supports, so each is zero.
Closedness makes the energy terms zero by the criterion following (B.4) of \cite{DNS}, hence $\alpha_j^2=0$. Nefness follows from its Corollary~2.4.
\end{proof}

\begin{corollary}\label{cor:domain}
Let $X\subset S$ be a continuously weakly complete noncompact domain in a compact K\"ahler surface. If $X$ is not a modification of a Stein space, every positive $\ddc$-closed current compactly supported in $X$ with zero Lelong numbers is closed. All such nonzero currents, normalized to mass one, represent one nef isotropic class $\alpha_X$.
\end{corollary}
\begin{proof}
Pair any given current with the current of Proposition~\ref{prop:averaging} in a sufficiently high disjoint shell. Proposition~\ref{prop:disjoint} applies. Compare two arbitrary normalized currents using a third current supported beyond both.
\end{proof}

\section{Why closedness does not finish the exhaustion problem}
\subsection*{A global $\partial\bar\partial$ lemma cannot be assumed}
There is an elementary failure even on a surface already possessing a smooth psh exhaustion. Let
$X=E\times\mathbb C$, where $E=\mathbb C/\Lambda$, and let $dw$ denote its invariant holomorphic form. With $z$ the coordinate on $\mathbb C$, put
\[
 \alpha=d\,\operatorname{Re}(z\,d\bar w)
        =\operatorname{Re}(dz\wedge d\bar w).
\]
This is a real, smooth, $d$-exact $(1,1)$-form. It is not $\D$-exact. If $\D u=\alpha$ distributionally, taking the trace and applying elliptic regularity makes $u$ smooth. Its restriction to each compact fiber $E\times\{z\}$ is harmonic and hence constant. Thus $u$ is independent of $w$, contradicting the nonzero mixed component of $\alpha$.
This example has noncompact support: it disproves the unrestricted replacement of the compact $\partial\bar\partial$ lemma, not a specially formulated statement for compactly supported shell currents.

\subsection*{Potentials of two shell currents need not separate levels}
Even in a compact ambient surface, the potential supplied by equality of classes need not distinguish entire exhaustion levels. Take $S=E\times\mathbb P^1$, $X=E\times\mathbb C$, and $\rho=|z|^2$.
Let $\mu_0,\mu_3$ be smooth radial probability measures supported in very small disks centered at $0$ and $3$, respectively. Their pullbacks $T_0,T_3$ are positive closed currents in ordered, disjoint $\rho$-bands, with smooth coefficients and identical cohomology classes.
Up to a positive factor and an additive constant, a potential $f$ for $T_0-T_3$ agrees outside those disks with
\[
 f(z)=\log|z|-\log|z-3|.
\]
The equality follows from the circular mean-value formula for the logarithmic kernel. Nevertheless,
\[
 \max_{|z|=1}f(z)=\log(1/2)
 \quad\text{and}\quad
 \min_{|z|=2}f(z)=\log(2/5)<\log(1/2).
\]
No single threshold separates these two complete levels in the expected order. This is not a counterexample to regularization; it shows why a choice of currents with additional separation properties is needed.

\subsection*{The remaining steps}
For an arbitrary $X$, Proposition~\ref{prop:averaging} supplies singular currents and the metric construction supplies completeness and infinite volume. A global curvature bound has not been obtained from these hypotheses, so the heat argument does not yet give general closedness. Nor would closedness alone supply a proper smooth psh function: a global potential or end-separation construction is still required.

One precise sufficient target is a cofinal nested sequence of compact sets $K_n$ and smooth psh functions $v_n\ge0$ on $X$ such that $v_n$ vanishes near $K_n$ and $v_n\ge1$ on $X\setminus K_{n+1}$. Then
$\sum_n v_n$ is locally finite, smooth, psh, and exhaustive. Indeed, outside $K_{N+1}$ its first $N$ summands are at least one. This criterion describes the missing construction. The following sections sharpen its dual formulation and establish it under an additional Green-kernel hypothesis.

Classifying continuous exhaustion levels as smooth Levi-flat hypersurfaces would also be unjustified. The relevant structural results in \cite[Theorem~1.1]{BM} require regular levels and $C^2$ or $C^4$ regularity; their introduction explicitly distinguishes the minimal kernels for different regularities.

The unrestricted implication remains unproved. The new separation results below give a conditional theorem and isolate the precise additional input needed for a general argument; they do not extend the manuscript's compact-divisor theorem to all K\"ahler surfaces.


\section{The exact separation problem and its duality gap}\label{sec:separation-duality}

% Begin separation-duality
The following statements isolate a precise remaining problem. Except where
K\"ahlerness is invoked for almost-psh regularization, they apply in any
complex dimension. Put
\[
 \mathcal P^\infty(X)=C^\infty(X,\mathbb R)\cap\operatorname{PSH}(X),
 \qquad
 \widehat K_\infty=
 \{p\in X:u(p)\le\sup_Ku\ \text{for every }u\in\mathcal P^\infty(X)\}.
\]
For a nonempty compact set $K$, let $J_p^\infty(K)$ consist of the
probability measures $\mu$ supported on $K$ such that
$u(p)\le\int_Ku\,d\mu$ for every $u\in\mathcal P^\infty(X)$.
These measures test the smooth cone only.

\begin{proposition}[Exact separator criterion]\label{prop:separator-criterion}
Let $K,L\subset X$ be disjoint nonempty compact sets. The following are
equivalent:
\begin{enumerate}
\item There is $u\in\mathcal P^\infty(X)$ with
$\sup_Ku<\inf_Lu$.
\item $L\cap\widehat K_\infty=\varnothing$.
\item $J_p^\infty(K)=\varnothing$ for every $p\in L$.
\end{enumerate}
\end{proposition}
\begin{proof}
If each $p\in L$ lies outside the hull, choose a smooth psh function,
after positive affine rescaling, smaller than $-2$ on $K$ and larger
than $3$ near $p$. Finitely many such neighborhoods cover $L$.
A regularized maximum of the corresponding functions, with sufficiently
small regularization parameter, is a separator. The converse is immediate.

For completeness, membership in the hull implies the existence of one
Jensen measure satisfying all the inequalities simultaneously. For any
finite list $u_1,\ldots,u_r\in\mathcal P^\infty(X)$, the compact convex set
\[
 \left\{\left(\int_Ku_i\,d\mu-u_i(p)\right)_{i=1}^r:
                   \mu\in\operatorname{Prob}(K)\right\}
\]
must meet $\mathbb R_+^r$. Otherwise finite-dimensional separation gives
nonnegative coefficients $c_i$, not all zero, with
$\sup_K\sum_i c_i u_i<\sum_i c_i u_i(p)$, contradicting the hull condition.
Compactness of $\operatorname{Prob}(K)$ and the finite intersection
property now give the required measure. The reverse implication follows
by bounding $\int_Ku\,d\mu$ by $\sup_Ku$.
\end{proof}

\begin{proposition}[From level separators to global barriers]
\label{prop:separator-barriers}
Let $\rho$ be a continuous psh exhaustion. Set $K_a=\{\rho\le a\}$ and
$L_b=\{\rho=b\}$, where $b>a$ and both sets are nonempty. A smooth psh
separator for $K_a,L_b$ produces a nonnegative smooth psh function $v$
which vanishes near $K_a$ and satisfies $v\ge1$ on $\{\rho\ge b\}$.
Consequently, $X$ has a smooth psh exhaustion if and only if, for every
sufficiently large $a$, some $b>a$ satisfies
$J_p^\infty(K_a)=\varnothing$ for every $p\in L_b$.
\end{proposition}
\begin{proof}
Normalize the separator so that $u<-2$ on $K_a$ and $u>3$ on $L_b$.
Let $M$ be a smooth regularized maximum which is at least the maximum
and agrees with it whenever the two arguments differ by more than $1$.
Define
\[
 w=\begin{cases}
      u,&\rho\le b,\\
      M(u,1),&\rho>b.
    \end{cases}
\]
The two formulas coincide on an open neighborhood of $L_b$ because
$u>3$ there. Thus $w$ is globally smooth and psh, despite the possible
nonsmoothness of $L_b$, and $w\ge1$ outside $\{\rho<b\}$.
Compose with a smooth convex nondecreasing function $\chi$ which is zero
on $(-\infty,0]$ and at least $1$ on $[1,\infty)$ to obtain $v$.

Apply this construction to $a_j\to\infty$. The resulting barriers $v_j$
vanish near $K_{a_j}$, so $\sum_jv_j$ is locally finite, smooth and psh.
For every $N$ it is at least $N$ outside the compact union of the first
$N$ corresponding sublevel sets; hence it is exhaustive. Conversely,
a smooth psh exhaustion separates $K_a$ from $L_b$ once $b$ is large
enough, by compactness of its sublevel sets.
\end{proof}

In particular, failure of smooth exhaustion forces some compact sublevel
$K_a$ to have $\widehat K_{a,\infty}$ meeting every level $L_b$, $b>a$.
If $J_p^0(K)$ denotes Jensen measures testing all continuous psh functions,
then $J_p^0(K)\subset J_p^\infty(K)$. For $p\in L_b$ and $K=K_a$,
the function $\rho$ proves $J_p^0(K)=\varnothing$. Replacing
$J_p^\infty$ by $J_p^0$ without further proof would assume the missing
regularization property.

\begin{theorem}[Bounded-mass Green-current duality]
\label{thm:bounded-mass-duality}
Fix a Hermitian form $\omega$, a compact set $F\supset K\cup\{p\}$,
and $m\ge0$. For $u\in C^\infty(X,\mathbb R)$ put
\[
 q_F(u)=\max\left\{0,-\inf_{x\in F}
              \lambda_{\min}(\D u;\omega)(x)\right\}.
\]
There exist $\mu\in\operatorname{Prob}(K)$ and a positive current $T$
of bidimension $(1,1)$ such that
\[
 \supp T\subset F,\qquad \langle T,\omega\rangle\le m,
 \qquad \D T=\mu-\delta_p
\]
if and only if
\begin{equation}\label{eq:quantitative-jensen}
 u(p)-\sup_Ku\le m q_F(u)
 \qquad\text{for every }u\in C^\infty(X,\mathbb R).
\end{equation}
Equivalently,
\[
 \sup\{u(p)-\sup_Ku:\D u\ge-\omega\text{ on }F\}\le m.
\]
\end{theorem}
\begin{proof}
Positive currents supported on $F$ with mass at most $m$ form a compact
convex set in the weak topology. Therefore
\[
 \mathcal A=\{\mu-\D T:
     \mu\in\operatorname{Prob}(K),\ T\ge0,
     \ \supp T\subset F,\ \langle T,\omega\rangle\le m\}
\]
is weakly compact and convex in the space of compactly supported
distributions. Its support function, evaluated at a smooth real $u$, is
\[
 \sup_{\sigma\in\mathcal A}\langle\sigma,u\rangle
       =\sup_Ku+m q_F(u).
\]
Indeed, the current term is maximized either by $T=0$, or by placing
mass $m$ at a point and a complex tangent line attaining the most
negative eigenvalue of $\D u$ on $F$. Hahn--Banach separation now gives
$\delta_p\in\mathcal A$ exactly when
\eqref{eq:quantitative-jensen} holds. The final formulation follows by
scaling when $q_F(u)>0$; if $q_F(u)=0$, scaling by arbitrarily large
constants forces $u(p)\le\sup_Ku$.
\end{proof}

\begin{proposition}[Continuous psh functions exclude exact Green currents]
\label{prop:continuous-green-exclusion}
Suppose $X$ is K\"ahler and $\rho\in C^0(X)\cap\operatorname{PSH}(X)$
satisfies $\rho(p)>\sup_K\rho$. Then no compactly supported positive
current satisfies $\D T=\mu-\delta_p$ with
$\mu\in\operatorname{Prob}(K)$.
\end{proposition}
\begin{proof}
By \cite[Corollary~9.3]{Demailly82} and Dini's theorem, on each fixed
compact set $F$ there are smooth $\rho_j$ converging uniformly to $\rho$
with $\D\rho_j\ge-\varepsilon_j\omega$ on $F$, where
$\varepsilon_j\to0$. If $F$ contains $\supp T\cup K\cup\{p\}$, then
\[
 \int_K\rho_j\,d\mu-\rho_j(p)
   =\langle T,\D\rho_j\rangle
   \ge-\varepsilon_j\langle T,\omega\rangle.
\]
Passing to the limit contradicts the assumed strict gap. Equivalently,
the supremum in Theorem~\ref{thm:bounded-mass-duality} is infinite:
divide suitable approximants by any positive Levi-error bounds tending
to zero.
\end{proof}

\paragraph{What unrestricted Hahn--Banach actually gives.}
Let $E=C^\infty(X,\mathbb R)$ and let $E'$ be its continuous dual.
Set
\[
 \mathcal C=\{\D T:T\ge0\text{ of bidimension }(1,1),
                                  \ \supp T\Subset X\}.
\]
The polar of $\mathcal C$ is exactly $\mathcal P^\infty(X)$: point masses
in individual complex tangent directions test positivity of the Levi
form. The bipolar theorem therefore gives only
\[
 (\mathcal P^\infty(X))^\circ
      =\overline{\mathcal C}^{\,\sigma(E',E)},
\]
where the polar means nonnegative pairing. Thus a smooth Jensen measure
gives $\mu-\delta_p$ in this closure, not necessarily in $\mathcal C$.
The closure may require nets; neither a common compact support nor a
mass bound for representing currents follows from this assertion.
Harvey--Lawson's exact Green-current theorem assumes strict convexity
\cite[Theorem~3.8]{HL}. Their polar calculation and closedness proof
\cite[Lemmas~3.10--3.11]{HL} separate these two steps explicitly.
The latter uses a strictly psh $f$ to estimate
$\varepsilon\langle T_j,\omega\rangle
 \le\langle\D T_j,f\rangle$ on a common compact support.
That coercive estimate is unavailable here.

Even if a sequence $T_j$ has a common compact support and
$\D T_j\to\mu-\delta_p$ distributionally, its masses must tend to
infinity by Proposition~\ref{prop:continuous-green-exclusion} and
compactness. Normalizing the masses yields, after passage to a
subsequence, a positive $\D$-closed current of mass one. Its Lelong
numbers need not vanish: these supports are not disjoint exhaustion
levels. Even a proof that this limiting current is $d$-closed would not
restore the lost nonzero source equation.

There is a further useful restriction. Such a sequence cannot have
$\D T_j$ equal to signed measures of uniformly bounded total variation.
Under that bound, distributional convergence on a common compact set
implies weak convergence against continuous functions. Choose the
almost-psh approximation $r_j$ of $\rho$ separately for each $T_j$, so
that $r_j\to\rho$ uniformly there and
\[
 \D r_j\ge-\eta_j\omega,
 \qquad \eta_j(1+\langle T_j,\omega\rangle)\longrightarrow0.
\]
Then $\langle\D T_j,r_j\rangle\ge-o(1)$, while the measure bound makes
the left side converge to $\int_K\rho\,d\mu-\rho(p)<0$.
The weak dual closure therefore loses essential control of the source
distributions as well as the current masses.

\paragraph{Why subtracting the limiting current is unjustified.}
The elementary positive-semidefinite matrix model
\[
 A_n=\begin{pmatrix}n&1\\1&1/n\end{pmatrix},\qquad
 L(A)=(A_{22},\operatorname{Re}A_{12})
\]
has $A_n\ge0$ and $L(A_n)\to(0,1)$, but $(0,1)$ is not in the image
of the positive cone: positivity and $A_{22}=0$ force $A_{12}=0$.
Meanwhile $A_n/n\to S=\operatorname{diag}(1,0)$ and $L(S)=0$.
No positive multiple of this kernel element can be subtracted:
$\det(A_n-cS)=-c/n<0$ for every $c>0$.
This is a functional-analytic illustration, not a geometric
counterexample. It shows why closedness of normalized Levi limits alone
does not prove closedness of the image cone. A remaining sufficient
geometric statement would be the restricted closed-range property
\[
 \overline{\mathcal C}\cap
   \{\mu-\delta_p:\mu\in\operatorname{Prob}(K)\}
 \subset\mathcal C
\]
for compact sublevels $K$ and points $p$ on suitable outer levels,
with closure in the topology specified above. Alternatively, one needs
a quantitative estimate such as \eqref{eq:quantitative-jensen} for the
smooth Jensen obstructions. Neither assertion is established here.

% End separation-duality


\section{What an end potential must accomplish}\label{sec:separation-envelope}

% Begin separation-envelope
\begin{lemma}[A pluriharmonic separator at infinity]
\label{lem:interval-separator}
Let $X$ be a connected noncompact complex manifold, of any dimension,
and let $K\Subset X$ be compact. Suppose $u:X\setminus K\to I$ is
pluriharmonic, where $I=(a,b)$ is an open interval whose endpoints may
be infinite. Assume that $u$ is proper at infinity into $I$: for every
compact interval $J\Subset I$, the set
\[
 \{x\in X\setminus K:u(x)\in J\}
\]
is relatively compact in $X$. Then $X$ admits a smooth psh exhaustion.
No extension of $u$ across $K$ is required.
\end{lemma}

\begin{proof}
Choose a smooth convex function $F:I\to\mathbb R$ tending to $+\infty$
at both endpoints of $I$. For example, take $t^2$ and subtract
$\log(t-a)$ if $a$ is finite and $\log(b-t)$ if $b$ is finite.
Thus $F$ has compact sublevel sets in $I$. The function $v=F(u)$
is smooth and psh on $X\setminus K$, because
\[
 \D v=F''(u)i\partial u\wedge\bar\partial u\geq0.
\]
It tends to $+\infty$ along every sequence escaping all compact
subsets of $X$, by the hypothesis on $u$.

Take a relatively compact smooth domain $D$ with $K\Subset D$.
A sufficiently small closed collar of $\partial D$ is disjoint from
$K$, and $v$ is bounded there. Choose $C$ larger than this bound,
and a smooth convex nondecreasing function $q:\mathbb R\to[0,\infty)$
which vanishes on $(-\infty,C]$ and equals $t$ plus a constant for
$t\geq C+1$. Define
\[
 \Phi(x)=
 \begin{cases}
  0,&x\in D,\\
  q(v(x)),&x\in X\setminus D.
 \end{cases}
\]
Both definitions are identically zero on a neighborhood of
$\partial D$. Hence $\Phi$ is globally smooth and psh, and its
sublevel sets are compact.
\end{proof}

In particular, if a global real distribution satisfies $\D u=T$
with $\supp T\Subset X$, then $u$ is smooth pluriharmonic off
$\supp T$ by elliptic regularity. Lemma~\ref{lem:interval-separator}
applies if the additional interval-separation condition holds.
Existence of a potential alone does not establish that condition.

\begin{remark}[Zero Monge--Amp\`ere measure in the residual case]
\label{rem:flat-envelope}
Let $B$ be a compact Riemann surface and $L\in\operatorname{Pic}^0(B)$
a non-torsion unitary flat line bundle. On $Y=L\setminus B$, put
$h=\log|v|$ in unitary flat fiber coordinates. Constant unitary
transition functions make $h$ a global smooth pluriharmonic function.
Each level of $h$ is compact, and $h:Y\to\mathbb R$ is proper.
The flat connection and a K\"ahler form on $B$ give a K\"ahler
metric on $Y$. Moreover $\mathcal O(Y)=\mathbb C$: the Laurent
coefficients of a holomorphic function are sections of $L^{-n}$,
which have no nonzero sections for $n\ne0$, while the zeroth
coefficient is constant on $B$.

The smooth psh exhaustion $\rho=h^2$ satisfies
\[
 (\D\rho)^2=0,
 \qquad \D\rho=2i\partial h\wedge\bar\partial h\ne0.
\]
More precisely, for $0<c<d$, on the relatively compact collar
$A=\{c<h<d\}$ the two functions
\[
 u_1=\frac{h^2-c^2}{d^2-c^2},
 \qquad u_2=\frac{h-c}{d-c}
\]
are smooth psh solutions of
\[
 (\D u)^2=0,\qquad u|_{\{h=c\}}=0,
 \qquad u|_{\{h=d\}}=1,
\]
but $u_1<u_2$ in $A$. Thus the homogeneous Monge--Amp\`ere
Dirichlet equation is not unique on these collars, even when
$\mathcal O(Y)=\mathbb C$; its vanishing does not imply
pluriharmonicity.

In this particular flat model the upper psh envelope is $u_2$.
Indeed, infinite unitary holonomy is dense in the circle, so every
leaf of the flat foliation is dense in its compact $h$-level.
A psh function on $A$ attains its maximum on such a level;
the maximum principle on the maximizing leaf and upper
semicontinuity make it constant on that whole level. It is
therefore a convex function of $h$, and lies below the affine
secant of its boundary bounds. This argument uses the flat
foliation and does not prove an envelope theorem for general
continuously exhausted K\"ahler surfaces.
\end{remark}

\begin{remark}[A bounded potential can occur for every normalization]
\label{rem:bounded-flat-potential}
In the same construction restrict to $X=\{a<h<b\}$ with finite
$a<b$. We still have $\mathcal O(X)=\mathbb C$, and
\[
 -\log(h-a)-\log(b-h)
\]
is a smooth psh exhaustion. Choose a smooth convex $\chi$ on
$\mathbb R$ with $\chi''\not\equiv0$ supported in a compact
subinterval of $(a,b)$. Then $T=\D\chi(h)$ is nonzero, smooth,
positive, closed, compactly supported, and has zero Lelong numbers.
Its global potential $\chi(h)$ is bounded. Every global
pluriharmonic function on $X$ is affine in $h$: the compact-level
maximum argument above makes it a function of $h$, and
pluriharmonicity makes that function affine. Consequently every
distributional global potential of $T$ is bounded, since the
difference of two potentials is pluriharmonic. Properness as a
map to all of $\mathbb R$ is therefore too strong a desired
normalization. The useful criterion is separation into an open
interval followed by convex blow-up at its endpoints, as in
Lemma~\ref{lem:interval-separator}.
\end{remark}

\begin{remark}[What annular envelopes and truncations do not supply]
Even an actual psh upper envelope need not be pluriharmonic in
dimension two. On $r<|z|<R$ in $\mathbb C^2$, the upper envelope
with constant inner and outer bounds $0,1$ is
\[
 u(z)=\frac{\log|z|-\log r}{\log R-\log r}.
\]
Restricting competitors to the complex lines through the origin
proves the envelope assertion by the subharmonic maximum
principle on planar annuli. Its Levi form has rank one everywhere:
$(\D u)^2=0$ but $\D u\ne0$. This is a counterexample to the
unqualified local inference, not a counterexample to a possible
additional rigidity theorem under $\mathcal O(X)=\mathbb C$.

Ordinary harmonic capacitary selection does not automatically
improve matters. The real harmonic function with the same
boundary values in this four-dimensional annulus is
\[
 H(z)=\frac{r^{-2}-|z|^{-2}}{r^{-2}-R^{-2}}.
\]
Its complex Hessian has one positive and one negative eigenvalue,
so it is not psh. Showing that a harmonic separator is psh would
indeed make it pluriharmonic, but that is an additional assertion.

Finally, even in the smooth case,
\[
 \D\chi(\rho)=\chi'(\rho)\D\rho
       +\chi''(\rho)i\partial\rho\wedge\bar\partial\rho.
\]
Support of $\chi''$ in a bounded interval does not make this
current supported on the corresponding shell: the first term
persists on the outer end. A nonconstant increasing convex
$\chi$ cannot become constant there. Cutting off the upper
slope instead generally destroys positivity. These operations
therefore do not furnish the required positive shell current
with a separating potential.
\end{remark}

% End separation-envelope


\section{Separation by a canonical Green potential}\label{sec:separation-green}
This section proves a conditional theorem without a compactification.
The metric and Green hypotheses are additional assumptions; they have
not been derived from a continuous psh exhaustion.
The use of scalar and Hodge heat to solve a Poincar\'e--Lelong equation
has a precedent in Ni--Tam \cite[\S2]{NiTam}. Here the surface trace
isometry and the vanishing primitive defect provide the required
large-time decay for the particular compactly supported current.

% Begin separation-green
\begin{theorem}\label{thm:green-separation}
Let $(X,\omega)$ be a connected complete K\"ahler surface of infinite
volume and uniformly bounded Riemannian curvature. Suppose its minimal positive scalar Green
kernel exists and is regular at infinity: for one (and hence every)
pole $o$,
\[
 G(x,o)\longrightarrow0\quad\text{as }x\to\infty.
\]
If there is a nonzero positive $(1,1)$-current $T$ with compact
support such that $\D T=0$ and $\nu(T,x)=0$ at every point, then
$X$ admits a smooth plurisubharmonic exhaustion.

More precisely, for $\mu=T\wedge\omega$, the negative Green
potential of $\mu$, with the normalization specified below, is a
global psh potential of $T$, is smooth and pluriharmonic off
$\supp T$, and tends uniformly to zero at infinity.
\end{theorem}

\paragraph{Conventions.}
Write $\D=i\partial\bar\partial$, $dV=\omega^2/2$, and let
$\Lambda$ be contraction with $\omega$. Let $\Delta_k$ denote the
nonnegative Hodge Laplacian on $k$-forms, $P_t=e^{-t\Delta_0}$, and
$H_t=e^{-t\Delta_2}$. Set $L=\Lambda\D$. There is a fixed
$\kappa>0$ such that
\begin{equation}\label{eq:normalization}
 \Delta_0=-\kappa L.
\end{equation}
For the usual compatible conventions $\kappa=2$, but retaining
$\kappa$ avoids dependence on this convention. The same K\"ahler
identity gives, for every closed real $(1,1)$-form $B$,
\begin{equation}\label{eq:closed-hodge}
 \Delta_2 B=-\kappa\D\Lambda B.
\end{equation}
This can also be checked locally by writing $B=\D b$ and using
commutation of $\D$ with the Laplacian.

\paragraph{The trace-lifting isometry and its intertwining property.}
The scalar $L^2$ kernel vanishes on a connected complete
infinite-volume manifold: the complete-metric integration-by-parts
theorem makes an $L^2$ harmonic function constant, and infinite
volume makes that constant zero. Essential self-adjointness of the
scalar Laplacian on $C_c^\infty$ therefore gives density of
$L(C_c^\infty)$ in real $L^2$.

For real $u\in C_c^\infty(X)$, the surface linear-algebra identity
and Stokes' theorem give
\begin{equation}\label{eq:isometry}
 \int_X(\D u)^2=0,
 \qquad \|\D u\|_2=\|Lu\|_2.
\end{equation}
Indeed, for real $(1,1)$-forms $B$,
$B^2=((\Lambda B)^2-|B|^2)dV$.
Consequently there is a well-defined linear isometry
\[
 R:L^2(X,\mathbb R)\longrightarrow
                   L^2(X,\Lambda^{1,1}_{\mathbb R})
\]
given by $Rf=\lim_j\D u_j$ whenever $Lu_j\to f$ in $L^2$.
It satisfies
\[
 d(Rf)=0,\qquad \Lambda Rf=f,
 \qquad \int_X(Rf)^2=0.
\]

For every $h>0$, this isometry intertwines heat:
\begin{equation}\label{eq:intertwining}
 H_h Rf=R P_h f.
\end{equation}
To justify this without assuming a bounded Green operator, first
take $f=Lu$ for $u\in C_c^\infty$. K\"ahler identities give
$H_h\D u=\D P_hu$, while $LP_hu=P_hLu$.
The function $P_hu$ lies in the operator domain of $L$.
Since $C_c^\infty$ is a core, approximate $P_hu$ in the graph norm
of $L$ by compactly supported smooth functions. Equation
\eqref{eq:isometry} then shows that
$\D P_hu=RLP_hu=RP_hLu$. Finally use density of $L(C_c^\infty)$
and boundedness of both heat operators and $R$ to obtain
\eqref{eq:intertwining} for all $f\in L^2$.
The complete-metric operator facts used here are in
\cite[\S12]{DemaillyHodge}.

For every real $L^2$ $(1,1)$-form $A$ satisfying $\D A=0$, one has
the primitive-defect identity
\begin{equation}\label{eq:defect}
 \beta:=A-R\Lambda A,
 \qquad
 \int_X A^2=-\|\beta\|_2^2.
\end{equation}
In fact, $\Lambda\beta=0$, and testing $\D A=0$ against the
functions approximating $\Lambda A$ gives
$\int_XA\wedge R\Lambda A=0$. Together with
$\int_X(R\Lambda A)^2=0$, this proves \eqref{eq:defect}.

\begin{proof}[Proof of Theorem~\ref{thm:green-separation}]
Put $A_t=H_tT$ and $q_t=P_t\mu$. Heat acts on the compactly
supported measure-valued form $T$ by its smooth kernel. K\"ahler
commutation, by transposition, gives
\[
 \D A_t=0,\qquad \Lambda A_t=q_t.
\]
The second identity means equality of functions after the scalar
distribution $\Lambda T$ is identified with $\mu$ relative to $dV$.

We recall the short-time estimate established in the accompanying
proof of Proposition~\ref{prop:bounded-geometry-closedness}. Vanishing Lelong numbers and
compact support give
$\sup_x\mu(B(x,r))=o(r^2)$ uniformly. Scalar Gaussian heat bounds
then imply
\[
 t\|q_t\|_2^2=o(1).
\]
The Weitzenb\"ock formula, connection-heat cone preservation, and
Duhamel's formula give
\[
 |A_t|\leq Cq_t,
 \qquad A_t\geq-Ctq_t\,\omega\qquad(0<t\leq1).
\]
Hence $A_t\in L^2$ and
$\int_XA_t^2\geq-Ct\|q_t\|_2^2=o(1)$.
For the underlying semigroup domination, Gaussian bounds, and
parallel-transport formula, see
\cite[Introduction and Lemmas~3.6--3.7]{BDG}; the full deduction for measure-valued $T$ appears in the proof of Proposition~\ref{prop:bounded-geometry-closedness}.

Define $\beta_t=A_t-Rq_t$. Equation~\eqref{eq:defect} implies
\begin{equation}\label{eq:small-defect}
 \|\beta_t\|_2\longrightarrow0\quad(t\downarrow0).
\end{equation}
The important additional observation is that this defect vanishes
at every positive time. If $0<s<t$, the semigroup identities and
\eqref{eq:intertwining} give
\[
 \beta_t=H_{t-s}A_s-RP_{t-s}q_s
         =H_{t-s}(A_s-Rq_s)=H_{t-s}\beta_s.
\]
The Hodge heat semigroup is an $L^2$ contraction. Letting $s\downarrow0$
in this equality and using \eqref{eq:small-defect}, we obtain
\begin{equation}\label{eq:exact-lift}
 A_t=Rq_t\quad(t>0).
\end{equation}
Thus $A_t$ is closed for every $t>0$.

Fix $s>0$. Since $q_s\in L^2$ and the scalar $L^2$ kernel is zero,
the spectral theorem gives $P_{t-s}q_s\to0$ in $L^2$ as $t\to\infty$.
Equation~\eqref{eq:exact-lift} and the isometry of $R$ imply
\begin{equation}\label{eq:large-time}
 A_t\longrightarrow0\quad\text{in }L^2\quad(t\to\infty).
\end{equation}

Because $A_t$ is closed, \eqref{eq:closed-hodge} gives
\[
 \frac{\partial A_t}{\partial t}
 =-\Delta_2A_t=\kappa\D q_t.
\]
Integration from $\epsilon$ to $b$ yields
\[
 A_b-A_\epsilon
      =\kappa\D\int_\epsilon^b P_t\mu\,dt.
\]
The positive Green kernel is $G(x,y)=\int_0^\infty p_t(x,y)\,dt$.
Its potential $G\mu$ is locally integrable: the local singularity
of $G$ is integrable against volume in the first variable, uniformly
for its second variable in a compact set, and $\mu$ has finite
compact support. Monotone convergence therefore applies in
$L^1_{\rm loc}$ to the displayed time integrals. Using
$A_\epsilon\to T$ as currents and \eqref{eq:large-time}, we conclude
\begin{equation}\label{eq:green-potential}
 T=\D v,\qquad v=-\kappa G\mu.
\end{equation}

Since $T\geq0$, the distribution $v$ has a psh representative.
It agrees with the negative Green potential; this potential is
upper semicontinuous, is strictly negative everywhere, and is not
identically $-\infty$. It is smooth off $K=\supp T$ and
pluriharmonic there by \eqref{eq:green-potential}.
Regularity of the Green kernel at infinity implies $v(x)\to0$
uniformly as $x\to\infty$. To see the required uniformity over
the source variable, fix a compact set containing $K$ and $o$ in
its interior. For $x$ outside a larger compact set, the Harnack
inequality in the source variable gives
$G(x,y)\leq C_KG(x,o)$ for all $y\in K$. Integrate against $\mu$.

Choose a relatively compact neighborhood $U$ of $K$ and a number
$c<0$ with $\sup_{\overline U}v<c$. Such a strict choice is
possible because $v$ is strictly negative and upper semicontinuous
on the compact set $\overline U$; possible values $-\infty$ cause
no problem. Take a smooth convex regularized maximum $M_\delta$
of two real variables, nondecreasing in each, equal to the larger
variable when their difference exceeds $\delta$.
Choose $\delta>0$ so small that
$v<c-2\delta$ near $\overline U$ and $c+2\delta<0$.
Then
\[
 w=M_\delta(v,c)
\]
is defined to be $c$ where $v=-\infty$; it is constant on $U$,
smooth and psh on all of $X$, and strictly negative. It equals
$v$ outside a sufficiently large compact set. Therefore
\[
 \Phi=-\log(-w)
\]
is a smooth psh exhaustion: the scalar function $s\mapsto-\log(-s)$
is convex and increasing for $s<0$, and $w\to0$ at infinity.
\end{proof}

\begin{remark}[The genuine remaining restriction]
The Green hypothesis is substantive. For example, $E\times\mathbb C$
with $E$ compact admits no nonconstant bounded-above psh function:
restriction to compact fibers makes such a function descend to a
bounded-above subharmonic function on $\mathbb C$, hence it is
constant. It therefore cannot satisfy all hypotheses of
Theorem~\ref{thm:green-separation}, which would produce the
nonconstant bounded-above psh function $v$ (or $w$). Nevertheless
it has a smooth psh exhaustion. The theorem accordingly covers a
regular hyperbolic situation and does not handle the parabolic
case by itself.

A continuous psh exhaustion can supply suitable hyperbolic
sublevel domains (compare \cite[proof of Theorem~2.5, p.~1360]{NR}),
but applying the theorem separately on such domains does not by
itself provide compatible global smooth separators on their union.
No such compatibility assertion is made here.

\end{remark}


\begin{remark}[Changing the positive source cannot repair an irregular Green end]
For any nonzero finite positive compactly supported measure $\mu$, Harnack's
inequality in the source variable gives positive constants $c_\mu,C_\mu$ and
a compact set $F$ such that
\[
 c_\mu G(x,o)\le G\mu(x)\le C_\mu G(x,o),\qquad x\notin F.
\]
Indeed, put $o$ and $\supp\mu$ in a connected relatively compact domain;
for $x$ outside a larger compact set, $y\mapsto G(x,y)$ is positive
harmonic on that domain, and a finite Harnack chain gives uniform
upper and lower comparisons. Integrating gives the displayed bounds.
Thus if $G(x_j,o)$ fails to tend to zero along an escaping sequence,
changing the compact positive source current cannot make its Green
potential tend to zero along that sequence. The Green construction
solves separation under a boundary condition at infinity; that condition
cannot be manufactured merely by reweighting averaged currents.
\end{remark}

\begin{remark}[Scalar Evans potentials do not supply the missing complex positivity]
Hansen--Netuka's theorem \cite{HN} supplies positive harmonic Evans
potentials tending to infinity outside nonpolar compact sets on
parabolic Riemannian manifolds. Such a potential need not be
pluriharmonic, even on a K\"ahler surface already admitting a smooth
psh exhaustion. For example, take the flat product $E\times\mathbb C$,
a real nonconstant eigenfunction $\Delta_E\varphi=\lambda\varphi$
with $\lambda>0$, and $0<\varepsilon\|\varphi\|_\infty<1$.
For $r=|z|>0$ put
\[
 u(w,z)=\log r+\varepsilon\varphi(w)K_0(\sqrt\lambda\,r),
\]
where $K_0$ is the decaying modified Bessel function. Its equation
$K_0''(s)+s^{-1}K_0'(s)-K_0(s)=0$ makes $u$ harmonic for the product
metric. Since $K_0'=-K_1$ and $sK_1(s)\le1$,
\[
 \partial_r u\ge \frac{1-\varepsilon\|\varphi\|_\infty}{r}>0.
\]
The asymptotics $K_0(s)=-\log s+O(1)$ at zero and $K_0(s)\to0$ at
infinity show that $\{u=0\}$ is a smooth compact radial graph, bounding
a compact neighborhood of $E\times\{0\}$, and that $u$ is positive and
divergent on its exterior. Nevertheless the restriction of $\D u$
to a compact $E$ fiber is
$\varepsilon K_0(\sqrt\lambda r)\D_E\varphi\not\equiv0$.
Thus $u$ is not pluriharmonic; being harmonic, it cannot be psh either.
What remains necessary is a pluriharmonic Evans potential, or an Evans
construction compatible with the complex current equation.
\end{remark}

% End separation-green


\begin{corollary}\label{cor:regular-green-smoothing}
Let $X$ be a connected noncompact K\"ahler surface with a continuous psh
exhaustion. If $X$ admits a complete infinite-volume K\"ahler metric
with uniformly bounded Riemannian curvature whose minimal positive
Green kernel tends to zero at infinity, then $X$ admits a smooth psh
exhaustion.
\end{corollary}
\begin{proof}
If $X$ is a modification of a Stein space, pull back a smooth strictly
psh exhaustion of that space. Otherwise Proposition~\ref{prop:averaging}
supplies a nonzero compactly supported positive $\D$-closed current
with zero Lelong numbers. Theorem~\ref{thm:green-separation} applies.
\end{proof}

The remaining general problem is to produce comparable end separation
without the regular Green hypothesis, while also removing the curvature
assumption from the singular-current argument. Neither an arbitrary
choice of averaged currents, a formal noncompact $\partial\bar\partial$
lemma, nor a homogeneous Monge--Amp\`ere equation supplies these steps.

\begin{thebibliography}{99}
\bibitem{DemaillyHodge}
J.-P. Demailly, \emph{$L^2$ Hodge Theory and Vanishing Theorems}, in \emph{Introduction to Hodge Theory}, SMF/AMS Texts and Monographs \textbf{8} (2002), \S12, Hodge theory of complete K\"ahler manifolds.
\href{https://www-fourier.univ-grenoble-alpes.fr/~demailly/manuscripts/hodge-ams.pdf}{Author-hosted text}.
\bibitem{Demailly82}
J.-P. Demailly, \emph{Estimations $L^2$ pour l'op\'erateur $\bar\partial$ d'un fibr\'e vectoriel holomorphe semi-positif au-dessus d'une vari\'et\'e k\"ahl\'erienne compl\`ete}, Ann. Sci. \'Ecole Norm. Sup. (4) \textbf{15} (1982), 457--511.
\href{https://www.numdam.org/item/ASENS_1982_4_15_3_457_0/}{Authoritative scan}. Theorems 0.8, 1.3, 9.4.
\bibitem{Dingoyan}
P. Dingoyan, \emph{A factorisation theorem for curves with vanishing self-intersection}, arXiv:1405.3593.
\href{https://arxiv.org/abs/1405.3593}{Preprint}, Lemma 3.1.1.
\bibitem{DNS}
T.-C. Dinh, V.-A. Nguy\^en, and N. Sibony, \emph{Unique ergodicity for foliations on compact K\"ahler surfaces}, Duke Math. J. \textbf{171} (2022), 2627--2698.
\href{https://doi.org/10.1215/00127094-2022-0044}{DOI};
\href{https://arxiv.org/pdf/1811.07450v2}{author version used for theorem numbering}.
\bibitem{MS}
S. Mongodi and Z. S{\l}odkowski, \emph{Domains with a continuous exhaustion in weakly complete surfaces}, Math. Z. \textbf{296} (2020), 1011--1019.
\href{https://doi.org/10.1007/s00209-020-02466-z}{DOI};
\href{https://arxiv.org/pdf/1904.03497}{preprint}.
\bibitem{Vajaitu}
V. V\^aj\^aitu, \emph{Holomorphic convexity of pseudoconvex surfaces}, Ann. Sc. Norm. Super. Pisa Cl. Sci. (5) \textbf{23} (2022), 1659--1678.
\href{https://doi.org/10.2422/2036-2145.202011_018}{DOI}.
\bibitem{NR}
T. Napier and M. Ramachandran, \emph{The Bochner--Hartogs dichotomy for weakly $1$-complete K\"ahler manifolds}, Ann. Inst. Fourier \textbf{47} (1997), 1345--1365.
\href{https://doi.org/10.5802/aif.1602}{DOI}.
\bibitem{BM}
F. Bianchi and S. Mongodi, \emph{On minimal kernels and Levi currents on weakly complete complex manifolds},
\href{https://arxiv.org/pdf/2102.05328}{arXiv:2102.05328}.
\bibitem{BDG}
R. Baumgarth, B. Devyver, and B. G\"uneysu,
\emph{Estimates for the covariant derivative of the heat semigroup on differential forms, and covariant Riesz transforms},
Math. Ann. \textbf{386} (2023), 1753--1798.
\href{https://doi.org/10.1007/s00208-022-02409-5}{DOI}.

\bibitem{HL}
F. R. Harvey and H. B. Lawson, Jr.,
\emph{Duality of positive currents and plurisubharmonic functions in calibrated geometry},
Amer. J. Math. \textbf{131} (2009), 1211--1239.
\href{https://doi.org/10.1353/ajm.0.0074}{DOI};
\href{https://www.math.stonybrook.edu/~blaine/phiPSH-Part2.pdf}{author version},
Theorem~3.8 and Lemmas~3.10--3.11.
\bibitem{HN}
W. Hansen and I. Netuka, \emph{On the existence of Evans potentials},
Math. Ann. \textbf{356} (2013), 1283--1302.
\href{https://doi.org/10.1007/s00208-012-0873-2}{DOI}.
\bibitem{NiTam}
L. Ni and L.-F. Tam, \emph{Poincar\'e--Lelong equation via the
Hodge--Laplace heat equation}, Compos. Math. \textbf{149} (2013), 1856--1870.
\href{https://doi.org/10.1112/S0010437X12000322}{DOI};
\href{https://mathweb.ucsd.edu/~lni/academic/PL-1109-v1.pdf}{author preprint}.
\end{thebibliography}
\end{document}
